\begin{figure}[ht]
\centering
\begin{tikzpicture}[x=1cm, y=1cm, font=\small]

% ============ 2-wire (left) ============
\begin{scope}[shift={(0, 0)}]
\node[anchor=west, font=\bfseries] at (-0.2, 4.6) {2-wire};

% current source symbol
\draw[thick] (0.6, 3.9) circle (0.28);
\node[font=\scriptsize] at (0.6, 3.9) {$I$};
\draw[thick] (0.6, 3.62) -- (0.6, 3.2);

% lead down (with lead resistance)
\draw[thick, engred] (0.6, 3.2) -- (0.6, 1.4);
\node[font=\scriptsize, color=engred, anchor=east] at (0.5, 2.3) {$R_L$};

% sensing element (RTD)
\draw[thick, engblue] (0.4, 1.4) rectangle (2.2, 0.9);
\draw[engblue, thick] (0.5, 1.35) -- (2.1, 0.95);
\node[font=\scriptsize, color=engblue, anchor=north] at (1.3, 0.85) {$R(T)$};

% return lead (with lead resistance)
\draw[thick, engred] (2.2, 1.4) -- (2.2, 3.2);
\node[font=\scriptsize, color=engred, anchor=west] at (2.3, 2.3) {$R_L$};
\draw[thick] (2.2, 3.2) -- (2.2, 3.62);

% voltmeter across the whole loop (top)
\draw[thick] (0.6, 3.9) -- (0.05, 3.9);
\draw[thick] (2.2, 3.9) -- (2.75, 3.9);
\draw[thick] (0.05, 3.9) -- (0.05, 3.2);
\draw[thick] (2.75, 3.9) -- (2.75, 3.2);
\draw[thick] (0.05, 3.2) -- (2.75, 3.2);
\node[font=\scriptsize] at (1.4, 3.05) {$V$};

\node[font=\scriptsize, anchor=north, align=center, color=engred]
  at (1.3, 0.4) {measures\\$R(T) + 2R_L$};
\end{scope}

% ============ 3-wire (middle) ============
\begin{scope}[shift={(5.0, 0)}]
\node[anchor=west, font=\bfseries] at (-0.2, 4.6) {3-wire};

\draw[thick] (0.6, 3.9) circle (0.28);
\node[font=\scriptsize] at (0.6, 3.9) {$I$};
\draw[thick] (0.6, 3.62) -- (0.6, 3.2);
\draw[thick, engred] (0.6, 3.2) -- (0.6, 1.4);
\node[font=\scriptsize, color=engred, anchor=east] at (0.5, 2.3) {$R_L$};

\draw[thick, engblue] (0.4, 1.4) rectangle (2.2, 0.9);
\draw[engblue, thick] (0.5, 1.35) -- (2.1, 0.95);
\node[font=\scriptsize, color=engblue, anchor=north] at (1.3, 0.85) {$R(T)$};

\draw[thick, engred] (2.2, 1.4) -- (2.2, 3.2);
\node[font=\scriptsize, color=engred, anchor=west] at (2.3, 2.3) {$R_L$};
\draw[thick] (2.2, 3.2) -- (2.2, 3.62);

% third (compensation) lead
\draw[thick, engamber] (1.3, 1.4) -- (1.3, 3.2);
\node[font=\scriptsize, color=engamber, anchor=west] at (1.35, 2.55) {$R_L$};
\draw[thick] (1.3, 3.2) -- (1.3, 3.62);
\node[font=\scriptsize, anchor=south] at (1.3, 3.62) {comp.};

\node[font=\scriptsize, anchor=north, align=center, color=engamber]
  at (1.3, 0.4) {bridge cancels\\one $R_L$};
\end{scope}

% ============ 4-wire (right) ============
\begin{scope}[shift={(10.0, 0)}]
\node[anchor=west, font=\bfseries] at (-0.2, 4.6) {4-wire (Kelvin)};

\draw[thick] (0.2, 3.9) circle (0.28);
\node[font=\scriptsize] at (0.2, 3.9) {$I$};
\draw[thick] (0.2, 3.62) -- (0.2, 3.2);

% force leads (outer)
\draw[thick, engred] (0.2, 3.2) -- (0.2, 1.4);
\draw[thick, engred] (2.6, 1.4) -- (2.6, 3.2);
\draw[thick] (2.6, 3.2) -- (2.6, 3.62);
\node[font=\scriptsize, color=engred, anchor=east] at (0.1, 2.3) {force};

\draw[thick, engblue] (0.2, 1.4) rectangle (2.6, 0.9);
\draw[engblue, thick] (0.35, 1.35) -- (2.45, 0.95);
\node[font=\scriptsize, color=engblue, anchor=north] at (1.4, 0.85) {$R(T)$};

% sense leads (inner, no current -> no drop)
\draw[thick, enggray] (0.8, 1.4) -- (0.8, 2.9);
\draw[thick, enggray] (2.0, 1.4) -- (2.0, 2.9);
\draw[thick, enggray] (0.8, 2.9) -- (0.4, 2.9);
\draw[thick, enggray] (2.0, 2.9) -- (2.4, 2.9);
\node[font=\scriptsize, color=enggray] at (1.4, 2.9) {$V$};

\node[font=\scriptsize, anchor=north, align=center, color=engblue]
  at (1.4, 0.4) {sense leads carry\\no current: $R_L$ vanishes};
\end{scope}

\end{tikzpicture}
\caption[RTD lead-wire connections: 2-, 3-, and 4-wire]{Lead
resistance $R_L$ in three RTD connection schemes. The 2-wire
connection adds $2R_L$ in series with the sensing element $R(T)$,
so a $1\,\Omega$ round-trip lead resistance appears as roughly
$2.6\,\text{K}$ of apparent temperature on a Pt100 at
$0\,\degC$. The 3-wire connection routes a compensation lead
that lets a balanced bridge subtract one $R_L$, cancelling the
lead resistance to first order when the two current-carrying
leads match. The 4-wire (Kelvin) connection forces the excitation
current through the outer leads and senses the voltage through the
inner leads, which carry no current: the lead drop vanishes to
the input impedance of the voltmeter, and only the true element
voltage is read.}
\label{fig:vol01:ch05:rtd-leadwire}
\end{figure}
